% !TeX program = xelatex
\documentclass[10pt,aspectratio=169]{beamer}

\usepackage{amsmath,amssymb,bm}
\usepackage{booktabs,array}
\usepackage{graphicx}
\usepackage{xcolor}
\usepackage{tikz}
\usepackage{etoolbox}
\usepackage{listings}

% ---- CJK (Overleaf 內建 Noto CJK 字型，xelatex 編譯) ----
\usepackage{xeCJK}
\setCJKmainfont{Noto Serif CJK TC}
\setCJKsansfont{Noto Sans CJK TC}
\setCJKmonofont{Noto Sans Mono CJK TC}

\definecolor{ntpublue}{RGB}{0,78,105}
\definecolor{ntpudark}{RGB}{0,55,75}
\definecolor{ntpulight}{RGB}{228,240,244}
\definecolor{accentred}{RGB}{177,36,44}
\definecolor{softgreen}{RGB}{226,242,235}

\setbeamertemplate{navigation symbols}{}
\setbeamercolor{normal text}{fg=black,bg=white}
\setbeamercolor{frametitle}{fg=ntpudark,bg=white}
\setbeamerfont{frametitle}{size=\large,series=\bfseries}
\setbeamerfont{title}{size=\Large}
\setbeamerfont{subtitle}{size=\normalsize}
\setbeamerfont{author}{size=\small}
\setbeamerfont{institute}{size=\small}
\setbeamerfont{date}{size=\small}
\setbeamertemplate{itemize item}{\color{ntpublue}\scriptsize$\blacksquare$}
\setbeamertemplate{itemize subitem}{\color{ntpublue}\scriptsize$\blacktriangleright$}

% ---- code style ----
\definecolor{codebg}{RGB}{248,248,248}
\definecolor{codeframe}{RGB}{200,200,200}
\definecolor{keyword}{RGB}{0,0,180}
\definecolor{comment}{RGB}{100,100,100}
\definecolor{string}{RGB}{0,140,0}
\lstdefinestyle{Pystyle}{
  language=Python,
  backgroundcolor=\color{codebg},
  frame=single,
  rulecolor=\color{codeframe},
  basicstyle=\scriptsize\ttfamily,
  keywordstyle=\color{keyword}\bfseries,
  commentstyle=\color{comment}\itshape,
  stringstyle=\color{string},
  breaklines=true,
  showstringspaces=false,
  tabsize=2
}
\lstset{style=Pystyle}

\newcommand{\runningsection}{}
\newcommand{\setrunningsection}[1]{\renewcommand{\runningsection}{#1}}
\newcommand{\papertitle}{相似投手推薦系統：K-Nearest Neighbors 配球特徵相似度分析}
\newcommand{\shortpapertitle}{相似投手推薦系統 -- KNN 相似度分析}
\newcommand{\presenter}{Shang-Ying Shiu}
\setbeamercolor{headerfooter}{fg=white,bg=ntpublue}

\setbeamertemplate{headline}{
  \leavevmode
  {\color{ntpublue}\rule{\paperwidth}{0.05cm}}
  \vspace{0.10cm}
}

\setbeamertemplate{footline}{
  \leavevmode
  \hbox{%
    \begin{beamercolorbox}[wd=.22\paperwidth,ht=0.38cm,dp=0.12cm,leftskip=0.20cm]{headerfooter}
      \fontsize{7.2}{8.0}\selectfont\presenter
    \end{beamercolorbox}%
    \begin{beamercolorbox}[wd=.56\paperwidth,ht=0.38cm,dp=0.12cm,center]{headerfooter}
      \fontsize{7.2}{8.0}\selectfont\shortpapertitle
    \end{beamercolorbox}%
    \begin{beamercolorbox}[wd=.22\paperwidth,ht=0.38cm,dp=0.12cm,leftskip=0pt plus 1fil,rightskip=0.20cm]{headerfooter}
      \fontsize{7.2}{8.0}\selectfont\insertframenumber
    \end{beamercolorbox}%
  }%
}

\newcommand{\tocrow}[2]{%
  \ifdefstring{\runningsection}{#2}%
    {\tikz[baseline=-0.62ex]\node[fill=ntpublue,text=white,font=\bfseries,
        minimum width=0.5cm,minimum height=0.5cm,inner sep=0pt]{#1};%
     \hspace{0.3cm}\textbf{\color{black}#2}}%
    {\tikz[baseline=-0.62ex]\node[fill=black!8,text=black!35,
        minimum width=0.5cm,minimum height=0.5cm,inner sep=0pt]{#1};%
     \hspace{0.3cm}\color{black!35}#2}%
  \par\vspace{0.3cm}%
}

\newcommand{\SectionToc}[1]{
  \setrunningsection{#1}
  \begin{frame}{Table of Contents}
    \vspace{0.2cm}
    \tocrow{1}{資料撈取與特徵工程}
    \tocrow{2}{標準化與 KNN 演算法}
    \tocrow{3}{參數設定}
    \tocrow{4}{輸出格式與小結}
  \end{frame}
}

\title{\papertitle}
\author{\presenter}
\date{\today}

\begin{document}

\setrunningsection{}
\begin{frame}
  \titlepage
  \vspace{-0.2cm}
\end{frame}

\SectionToc{資料撈取與特徵工程}

\begin{frame}{系統目標與整體流程}
  \begin{itemize}
    \item \textbf{目標}：給定一位投手，從資料庫中所有投手裡找出配球特徵
          最相似的 5 位投手，供 App 端做「相似投手推薦」使用。
    \item 整體流程：
    \begin{itemize}
      \item 從 Supabase（\texttt{pitches} 資料表）撈取逐球資料，依「投手
            $\times$ 球種」聚合出配球比例與物理特徵
      \item 將每位投手整理成一個高維特徵向量，並做標準化
      \item 以 KNN 計算投手間的歐氏距離，取出最相似的 5 位投手
      \item 同時輸出成 Excel 報表與寫回 Supabase 資料表
    \end{itemize}
  \end{itemize}
\end{frame}

\begin{frame}[fragile]{原始資料聚合}
  以 SQL 依 \texttt{player\_name} 與 \texttt{pitch\_type} 分組，計算每位
  投手每種球路的物理特徵：
  \begin{lstlisting}
SELECT player_name, pitch_type,
       COUNT(*) as pitch_count,
       AVG(release_speed) as avg_speed,
       AVG(pfx_x) as avg_pfx_x,
       AVG(pfx_z) as avg_pfx_z,
       AVG(release_spin_rate) as avg_spin_rate
FROM pitches
WHERE pitch_type IS NOT NULL
GROUP BY player_name, pitch_type
  \end{lstlisting}
  \vspace{0.1cm}
  平均球速、平均水平位移（\texttt{pfx\_x}）、平均垂直位移
  （\texttt{pfx\_z}）、平均轉速，以及球路投出次數。
\end{frame}

\begin{frame}{樣本數過濾與球種使用率}
  \textbf{樣本數過濾}：只保留生涯總投球數足夠的投手，避免樣本過少導致
  統計量不穩定：
  \[
    \text{total\_pitches}(p)=\sum_{\text{pitch\_type}}\text{pitch\_count}(p,\text{pitch\_type})
  \]
  \[
    \text{保留條件：}\quad \text{total\_pitches}(p)\ge \texttt{min\_total\_pitches}=200
  \]
  \vspace{0.2cm}
  \textbf{球種使用率}：
  \[
    \text{usage\_rate}(p,t)=\frac{\text{pitch\_count}(p,t)}{\text{total\_pitches}(p)}
  \]
\end{frame}

\begin{frame}[fragile]{特徵矩陣建構}
  用 \texttt{pivot} 把「長表」（每列為 投手+球種）轉成「寬表」，每一列
  代表一位投手，欄位命名為 \texttt{特徵\_球種}：
  \begin{lstlisting}
df_pivot = df_model.pivot(index='player_name',
                           columns='pitch_type')
df_pivot.columns = [f"{col[0]}_{col[1]}"
                     for col in df_pivot.columns]
df_pivot = df_pivot.fillna(0)
  \end{lstlisting}
  \begin{itemize}
    \item 投手沒投過的球種，對應欄位以 0 填補
    \item 最終每位投手是一個高維向量，涵蓋各球種的
          \textbf{使用率、球速、轉速、水平位移、垂直位移}
  \end{itemize}
\end{frame}

\SectionToc{標準化與 KNN 演算法}

\begin{frame}{特徵標準化}
  各特徵單位量綱差異極大（使用率 0--1，球速 90+ mph，轉速 2000+ rpm）。
  若不標準化，轉速會在歐氏距離中完全主導計算，掩蓋使用率與位移的影響。
  \vspace{0.2cm}

  以 \texttt{StandardScaler} 做 Z-score 標準化：
  \[
    z_{ij}=\frac{x_{ij}-\mu_j}{\sigma_j}
  \]
  其中 $x_{ij}$ 為第 $i$ 位投手第 $j$ 個特徵的原始值，$\mu_j,\sigma_j$
  為該特徵在所有投手中的平均與標準差。
\end{frame}

\begin{frame}[fragile]{最近鄰搜尋}
  使用 \texttt{sklearn.neighbors.NearestNeighbors}：
  \begin{lstlisting}
scaler = StandardScaler()
X_scaled = scaler.fit_transform(X)

knn = NearestNeighbors(n_neighbors=n_neighbors + 1,
                        algorithm='auto')
knn.fit(X_scaled)
distances, indices = knn.kneighbors(X_scaled)
  \end{lstlisting}
  距離度量採用預設值，即\textbf{歐氏距離}：
  \[
    d(p,q)=\sqrt{\sum_{j=1}^{D}(z_{pj}-z_{qj})^2}
  \]
  距離值愈小代表兩位投手的配球特徵愈相似。
\end{frame}

\begin{frame}[fragile]{排除投手自己}
  每位投手在向量空間中距離自己最近的一定是自己本身（距離 $=0$），因此
  查詢時多取一個鄰居（\texttt{n\_neighbors + 1}），再從排名 1 開始取值：
  \begin{lstlisting}
for rank in range(1, n_neighbors + 1):
    sim_idx = indices[i][rank]
    sim_pitcher = pitcher_names[sim_idx]
    dist = distances[i][rank]
  \end{lstlisting}
  \begin{itemize}
    \item 排名 0：投手自己（距離 $=0$，捨棄不用）
    \item 排名 1--5：真正的 5 位相似投手
  \end{itemize}
\end{frame}

\SectionToc{參數設定}

\begin{frame}{參數設定總表}
  \begingroup\scriptsize
  \begin{tabular}{lll}
    \toprule
    參數 & 設定值 & 設定原因 \\
    \midrule
    \texttt{min\_total\_pitches} & 200 & 過濾投球樣本數過少的投手 \\
    \texttt{n\_neighbors} & 5 & 每位投手輸出 5 位最相似投手 \\
    \texttt{max\_pitches\_to\_show} & 5 & Excel 中每位投手最多列前 5 種球路 \\
    球種篩選門檻 & usage $>0$ & 只排除完全沒投過的球種，其餘皆視為真實球路 \\
    標準化方式 & StandardScaler & 消除特徵量綱差異，使距離計算公平 \\
    KNN 演算法 & \texttt{algorithm='auto'} & 依資料量/維度自動選最有效率的搜尋法 \\
    距離度量 & 歐氏距離（預設） & 值愈小代表兩位投手愈相似 \\
    \bottomrule
  \end{tabular}
  \endgroup
\end{frame}

\SectionToc{輸出格式與小結}

\begin{frame}{輸出格式}
  同一次運算同時產出兩種格式的結果：
  \vspace{0.2cm}

  \textbf{Excel 報表}（\texttt{pitcher\_similarities.xlsx}，人工檢視用）
  \begin{itemize}
    \item 每位目標投手先列出自己的基準列（$\star$ 基準，排名 0）
    \item 接著列出 5 位相似投手（排名 1--5），含前 5 種球路的使用率/
          球速/轉速/位移明細
  \end{itemize}

  \vspace{0.2cm}
  \textbf{Supabase 資料表}（\texttt{pitcher\_similarities}，API 查詢用）
  \begin{itemize}
    \item 欄位精簡為六項：\texttt{target\_pitcher}、\texttt{similar\_pitcher}、
          \texttt{rank}、\texttt{distance\_score}、\texttt{target\_stats}
          (JSONB)、\texttt{similar\_stats} (JSONB)
    \item 每次執行前先 \texttt{DROP TABLE IF EXISTS} 重建，確保資料與最新
          運算結果一致
  \end{itemize}
\end{frame}

\begin{frame}{小結}
  \begin{itemize}
    \item 以「球種使用率 + 球速 + 轉速 + 水平/垂直位移」構成每位投手的
          配球特徵向量
    \item 標準化後以 KNN（歐氏距離）計算投手間相似度
    \item 用最小樣本數門檻（200 球）過濾投球紀錄過少的投手
    \item 同時輸出人工可讀的 Excel 報表與供後端 API 查詢的資料庫表，
          兩者共用同一份運算結果，確保呈現一致
  \end{itemize}
\end{frame}

\setrunningsection{}
\begin{frame}
  \centering
  \vfill
  {\Large Thank you.}
  \vfill
\end{frame}

\end{document}
